% 问题四
\section{问题四：混合干扰源的完备发现与清除}

\subsection{模型准备}

\subsubsection{定向可见性与检测结果语义}

设源 $G$ 的定向发射方向为单位向量 $\boldsymbol o$，即 $\boldsymbol o$ 从源指向其有效覆盖区域的中心方向。对测点 $S$，检测点相对源的位置向量为 $S-G$。定向源的可见条件为

\begin{equation}
\boldsymbol o^{\mathsf T}(S-G)\ge0.
\label{eq:q4_visibility}
\end{equation}

该不等式表示测点落在以定向方向为中心的闭半平面内，恰好对应“定向方向两侧各 $90^\circ$ 且包含边界”的题面规则。于是检测结果可以统一写为

\begin{equation}
\operatorname{result}(S,c)=
\begin{cases}
\texttt{near},&
\lVert G-S\rVert_2\le5,\quad \boldsymbol o^{\mathsf T}(S-G)\ge0,\\[2mm]
\texttt{direction},&
5<\lVert G-S\rVert_2\le r^{\mathrm{recv}},\quad
\boldsymbol o^{\mathsf T}(S-G)\ge0,\\[2mm]
\texttt{no\_signal},&
\text{其他情形}.
\end{cases}
\label{eq:q4_measure_semantics}
\end{equation}

其中 \texttt{direction} 还满足示向误差约束

\begin{equation}
\left|\operatorname{wrap}_{180}(\hat\theta_c(S)-\theta(S,G))\right|\le\delta,
\label{eq:q4_bearing_error}
\end{equation}

$\operatorname{wrap}_{180}(\cdot)$ 将角差化为 $[-180^\circ,180^\circ)$ 内的有符号角。上式说明：\texttt{no\_signal} 可能由三种原因产生，即频道不存在、距离超过接收半径，或测点位于定向源背面。问题四不得把第三种原因误判为前两种。

\subsubsection{正证据更新与状态机}

当频道 $c$ 在某测点获得正常示向度时，先把该频道的可行区域多边形外包围 $A_c$ 初始化为目标圆域的外逼近多边形，再依次与示向锥 $W(S,\hat\theta_c)$ 求交，并保留接收距离不超过 $1500\mathrm{m}$ 的部分。若得到多个正常示向度，则继续执行凸多边形裁剪\cite{sutherland1974}。所有裁剪均使用 $1.01^\circ$ 安全半角，保证真实可行集始终包含在数值区域中。

由于 \texttt{no\_signal} 在问题四中具有方向歧义，问题三中“在已成功观测点与无信号点之间比较距离”的阴性半平面推理不能直接沿用。问题四只在定向源分支使用正常示向度作为位置约束，不在未知定向方向下用 \texttt{no\_signal} 收缩可行区域。该处理牺牲了部分可利用信息，但避免了由盲区产生的错误排除。

每个频道维护六类状态：未发现的 \texttt{unknown}，定位中的 \texttt{localizing}，待清除的 \texttt{clear\_ready}，已清除的 \texttt{cleared}，确认不存在的 \texttt{absent\_certified}，以及几何故障 \texttt{geometry\_fault}。频道可在以下三种情形下结束：清除指令返回成功；该频道在全部扫描站均无正响应，由方向完备扫描证书证明无源；或已发现的不同频道数达到 16，此时由题设 $N_J\le16$ 直接判定其余频道不存在。只有 20 个频道全部进入已清除或确认不存在状态时，任务才满足终止条件。

\subsubsection{时间模型与优化目标}

设路径依次访问位置 $P_0,P_1,\ldots,P_M$，则问题四的虚拟时间为

\begin{equation}
T=\sum_{k=1}^{M}\frac{\lVert P_k-P_{k-1}\rVert_2}{5}
+5N_{\mathrm{measure}}+N_{\mathrm{switch}}
+5N_{\mathrm{clear}}^{\mathrm{success}}+3N_{\mathrm{clear}}^{\mathrm{miss}}.
\label{eq:q4_total_time}
\end{equation}

式中各动作严格串行，清除动作不改变当前测向频道。

由于单次检测、换频和清除的耗时固定，而移动时间与总路径长度成正比，在线优化的主要目标是减少重复往返。每完成一个联合节点后，算法删除已闭合的扫描站或服务点，更新各频道的可行区域，并重新求解当前节点集上的开放路径。该动态重规划不改变扫描点集，因此不会削弱发现证书。

\subsection{模型建立}

\subsubsection{方向完备发现判据}

对固定位置 $G$，记位于接收半径下界内的全部扫描点相对向量为

\begin{equation}
V_G=
\left\{
P_j-G:
\ P_j\in P_{\mathrm{scan}},\ 
\lVert P_j-G\rVert_2\le1000
\right\},
\label{eq:q4_relative_set}
\end{equation}

其中 $P_{\mathrm{scan}}$ 可以为基准扫描集 $P_4$ 或提交版本的 $P_4'$。
由于真实接收半径满足 $r^{\mathrm{recv}}\ge1000\mathrm{m}$，$V_G$ 中每个扫描点都必定位于真实接收范围内。

\textbf{定理 1（方向完备发现判据）}　若对目标圆域内每个可能位置 $G$，均满足

\begin{equation}
0\in\operatorname{conv}V_G,
\label{eq:q4_direction_certificate}
\end{equation}

则无论定向方向如何，源 $G$ 都至少会在一个扫描点产生有效响应。反之，若某个 $G$ 不满足上式，则存在一个定向方向和接收半径 $r^{\mathrm{recv}}=1000\mathrm{m}$，使该源对所有扫描点均返回 \texttt{no\_signal}。

证明：先证充分性。设 $0\in\operatorname{conv}V_G$，任取单位发射方向 $\boldsymbol o$。若所有 $v\in V_G$ 都满足 $\boldsymbol o^{\mathsf T}v<0$，则因 $0$ 可写成 $V_G$ 的凸组合，取组合系数 $\lambda_v\ge0$（$\sum_{v\in V_G}\lambda_v=1$），有
\[
0=\boldsymbol o^{\mathsf T}0=\sum_{v\in V_G}\lambda_v\,\boldsymbol o^{\mathsf T}v<0,
\]

产生矛盾。因此至少存在一个 $v=P-G\in V_G$，使

\[
\boldsymbol o^{\mathsf T}(P-G)\ge0,
\qquad
\lVert P-G\rVert_2\le1000\le r^{\mathrm{recv}}.
\]

该扫描点位于定向源正面且处于有效接收范围内，故必返回 \texttt{direction} 或 \texttt{near}。

再证必要性。若 $V_G=\varnothing$，则接收半径下界内不存在可用扫描点，该源对任意方向都不可见，结论成立。以下设 $V_G\ne\varnothing$：若 $0\notin\operatorname{conv}V_G$，由有限点集凸包的闭性与严格分离定理，存在单位向量 $\boldsymbol o$ 和常数 $\alpha<0$，使

\[
\boldsymbol o^{\mathsf T}v\le\alpha<0,
\qquad
\forall v\in V_G.
\]

取该方向为定向源方向，并令 $r^{\mathrm{recv}}=1000\mathrm{m}$。此时所有位于接收半径内的扫描点都严格位于盲区，半径外的扫描点也不可能收到信号，因此全部测量均为 \texttt{no\_signal}。定理得证。

扫描集构造只依赖目标圆域和接收半径下界。
完整扫描后仍无正响应的频道可认证为无源；累计发现 16 个不同频道后，由 $N_J\le16$ 可直接将其余频道标记为确认不存在。

\subsubsection{基准 31 点等边三角格}

为得到证明简明、可用于对照的基准构造，取边长 $s=990\mathrm{m}$ 的等边三角格

\begin{equation}
e_1=(s,0),\qquad
e_2=\left(\frac{s}{2},\frac{\sqrt{3}s}{2}\right),
\qquad
\Lambda_s=
\left\{
ie_1+je_2:\ i,j\in\mathbb Z
\right\}.
\label{eq:q4_lattice}
\end{equation}

扫描集取与目标圆域相交的基本三角形的顶点并集

\begin{equation}
P_4=
\left\{
P\in\Lambda_s:
\lVert P\rVert_2\le1800+s
\right\}.
\label{eq:q4_base_sites}
\end{equation}

在式~\eqref{eq:q4_lattice} 给定的格点相位下，式~\eqref{eq:q4_base_sites} 枚举得到 31 个扫描点。对任意 $G\in\mathcal B_{1800}$，存在一个基本三角形包含 $G$，其三顶点均在 $P_4$ 中，到 $G$ 的距离不超过 $s=990\mathrm{m}$，且 $G$ 是三个顶点的凸组合。因此

\begin{equation}
0\in\operatorname{conv}V_G,
\label{eq:q4_lattice_certificate}
\end{equation}

该构造对任意朝向均至少可见，且证明直接，仅作为可核验基准，不主张其点数全局最少。

\subsubsection{22 点连续域认证扫描集}

提交版本进一步采用 22 点紧凑扫描集。记

\begin{equation}
P_4'=
\{(0,0)\}
\cup
\left\{
975d\left(\frac{360i}{7}\right):
i=0,\ldots,6
\right\}
\cup
\left\{
1850d\left(\frac{360j}{14}\right):
j=0,\ldots,13
\right\}.
\label{eq:q4_compact_sites}
\end{equation}

该集合由原点、半径 $975\mathrm{m}$ 的内环 7 点和半径 $1850\mathrm{m}$ 的外环 14 点组成，合计 22 点，其空间结构如图~\ref{fig:q4_compact_sites}所示。

\begin{figure}[htbp]
\centering
\begin{tikzpicture}[scale=1.45]
  \fill[black!3] (0,0) circle (1.8);
  \draw[thick] (0,0) circle (1.8);
  \draw[blue!45,dashed] (0,0) circle (0.975);
  \draw[red!45,dashed] (0,0) circle (1.85);
  \fill[black] (0,0) circle (2pt);
  \foreach \k in {0,...,6}{
    \fill[blue!75!black] ({0.975*cos(360*\k/7)},{0.975*sin(360*\k/7)}) circle (2pt);
  }
  \foreach \k in {0,...,13}{
    \fill[red!75!black] ({1.85*cos(360*\k/14)},{1.85*sin(360*\k/14)}) circle (2pt);
  }
  \coordinate (G) at (1.8,0);
  \draw[green!55!black,densely dashed] (G) circle (1.0);
  \fill[green!45!black] (G) circle (2.2pt);
  \node[below left] at (0,0) {原点};
  \node[above] at (0,1.8) {$\mathcal B_{1800}$};
  \node[right] at (1.8,0) {$G$};
  \node[green!45!black] at (2.35,0.65) {$1000\,\mathrm m$ 邻域};
\end{tikzpicture}
\caption{22 点紧凑扫描集及一个边界位置的接收邻域示意}
\label{fig:q4_compact_sites}
\end{figure}

由于 22 点集合不是完整三角格，其方向完备性采用连续域递归证书，而不是随机采样验证。
首先用目标圆域的外接正 72 边形覆盖 $\mathcal B_{1800}$，并划分为 72 个三角形。
对任一三角形单元 $\tau$，定义可用扫描点集合

\begin{equation}
E(\tau)=
\left\{
P\in P_4':
\max_{X\in\operatorname{vert}(\tau)}
\lVert P-X\rVert_2\le994.99
\right\}.
\label{eq:q4_eligible_set}
\end{equation}

若 $\tau\subseteq\operatorname{conv}E(\tau)$，则该单元通过认证。事实上，三个顶点都在凸包内意味着整个三角形都在凸包内；又因欧氏范数是凸函数，距离在每个三角形上的最大值必在某个顶点取得，所以 $\tau$ 中任意点 $G$ 到任意 $P\in E(\tau)$ 的距离均不超过 $994.99\mathrm{m}<1000\mathrm{m}$，从而

\begin{equation}
0\in\operatorname{conv}\{P-G:P\in E(\tau)\}\subseteq\operatorname{conv}V_G.
\label{eq:q4_cell_certificate}
\end{equation}

若单元尚未通过，则沿最长边递归二分；程序设置最大深度 22 和最多 60000 个单元。
程序核验表明，该集合在 7372 个单元内完成认证；一旦达到限制仍未完成，程序拒绝该证书并转入 25 点双环工程回退分支。该点集由原点、半径 $980\,\mathrm m$ 和 $1870\,\mathrm m$ 的两个正十二边形构成；启用前检查其扇形三角剖分的最大边长小于 $999\,\mathrm m$，且外环内切圆半径满足 $1870\cos15^\circ>1800\,\mathrm m$。前者保证三角形内任意源到三个顶点的距离均小于 $1000\,\mathrm m$，后者保证剖分覆盖目标圆域，故定理 1 的条件成立；任一不等式不满足时停止执行。该回退分支在本文 18 局完全提速版演练和三次正式测试中均未触发。

\subsubsection{定向源的定位与正证据收缩}

对频道 $c$ 的首次正常示向度，初始化

\begin{equation}
K_c=
\mathcal B_{1800}
\cap
W(S_0,\hat\theta_c)
\cap
\left\{
G:
5<\lVert G-S_0\rVert_2\le1500
\right\}.
\label{eq:q4_initial_region}
\end{equation}

数值实现使用覆盖 $K_c$ 的外逼近多边形 $A_c$，并始终保持

\begin{equation}
K_c\subseteq A_c.
\label{eq:q4_outer_contain}
\end{equation}

每次获得新的正常示向度后，执行

\begin{equation}
A_c\leftarrow
A_c\cap W(S,\hat\theta_c)\cap
\left\{
G:
\lVert G-S\rVert_2\le1500
\right\}.
\label{eq:q4_region_update}
\end{equation}

设 $C_c$ 为 $A_c$ 的覆盖圆心，定义保守半径

\begin{equation}
\rho_c=\max_{X\in A_c}\lVert X-C_c\rVert_2.
\label{eq:q4_outer_radius}
\end{equation}

由于欧氏范数是凸函数，最大值在凸多边形顶点取得。若

\begin{equation}
\rho_c\le r_{\mathrm{opt}}-\epsilon_{\mathrm{num}}
=19.999\mathrm{m},
\label{eq:q4_clear_certificate}
\end{equation}

则可行集内任意真实源到 $C_c$ 的距离都小于 $20\mathrm{m}$；由于清除成功只取决于该距离，在 $C_c$ 处执行清除即可通过距离证书保证成功。

若继续获得的都是正常示向度，上述正证据收缩与问题三一致。区别在于：当后续测量返回 \texttt{no\_signal} 时，问题四不将其视为几何矛盾，而是转入覆盖式清除分支。
原因是首测点可能恰位于定向源覆盖边界，且源距等于接收半径下界；
此时对任一备选下一测站 $S\ne S_0$，都存在与首测读数相容的源位置和定向方向，使 $S$ 落入背面；
因此不能保证下一测站必然收到信号，也不能将后续 \texttt{no\_signal} 解释为几何矛盾。

\subsubsection{与朝向无关的覆盖式清除}

设首次正常示向度为 $\hat\theta_0$，在线安全半角为 $\delta_{\mathrm{num}}$。首次观测成功说明真实源位于扇形

\begin{equation}
\mathcal R_0=
\left\{
G=S_0+\rho d(\phi):
5<\rho\le1500,\ 
\left|\operatorname{wrap}_{180}(\phi-\hat\theta_0)\right|
\le\delta_{\mathrm{num}}
\right\}.
\label{eq:q4_first_sector}
\end{equation}

在 $\mathcal R_0$ 上铺设等边三角格，并只保留到 $\mathcal R_0$ 的距离不超过 $\rho_{\mathrm{cov}}$ 的格点。采用边长

\begin{equation}
s_c=20\sqrt{3}<35\mathrm{m}
\label{eq:q4_clear_spacing}
\end{equation}

时，平面等边三角格对任意点的覆盖半径不超过 $s_c/\sqrt3=20\mathrm{m}$。数值实现将边长再乘以 $0.995$ 的安全系数，得到

\begin{equation}
\rho_{\mathrm{cov}}=19.9\mathrm{m}.
\label{eq:q4_clear_cover_radius}
\end{equation}

因此，对任意 $G\in\mathcal R_0$，存在一个格点 $Q$ 满足

\begin{equation}
\lVert Q-G\rVert_2\le19.9\mathrm{m}<20\mathrm{m}.
\label{eq:q4_clear_guarantee}
\end{equation}

清除是否成功与朝向无关，故有限格点序列中至少一次成功。
该过程不依赖后续测向，以有限次移动和清除换取方向盲区下的完备性。

提交版本按证明任务使用三个半径：理论覆盖半径不超过 $20\mathrm{m}$，数值实现按 $0.995$ 安全系数取 $19.9\mathrm{m}$；预算内分支以 $19.5\mathrm{m}$ 单元中心覆盖当前外包围 $A_c$，余下预算足以遍历全部单元时必命中。
$19.5\mathrm{m}$ 仅用于预算遍历证书，实际清除仍以接口距离检查确认。
预算不足时，改用 12 点近距测站包收缩区域；进入最后单元时，要求该单元全部顶点到历史失败圆或当前站圆心的最远距离不超过 $19.999\mathrm{m}$。真源不在失败圆内，故当前站圆必包含真源。
因此，每局 5 次空清除只是代价控制，不是完成性前提。

\subsubsection{发现、定位与清除的完备性}

在总扫描阶段，若累计发现的不同频道数达到 16，则由 $N_J\le16$ 可知其余频道均无源；否则扫描继续执行。对任一实际存在且尚未被清除的干扰源，由定理 1 和 22 点连续域证书可知，至少有一个扫描点能够看到它。因此扫描不会遗漏仍可能存在的源，也不可能把“源存在但所有扫描点均无信号”的频道错误认证为空。

对已发现频道，若正常示向度持续给出并使 $\rho_c\le19.999\mathrm{m}$，则由外包围包含关系直接得到清除距离证书。
出现 \texttt{no\_signal} 时，前述覆盖证书已保证有限次清除内至少成功一次；预算、12 点测站包和最后单元证明仅优化成本，不削弱该保证。
若出现空区域、认证清除失败或覆盖证明无法闭合，算法进入 \texttt{geometry\_fault}，不会把异常改写为“确认不存在”。

因此，终止时每个频道都已清除，或由方向完备证书或 16 个已发现频道的上界证明无源，不会留下 \texttt{unknown} 或 \texttt{localizing} 状态。

\subsubsection{动态路径调度}

设当前重规划时刻仍需访问的扫描站和待服务频道服务点组成节点集 $V_t$，其节点数为 $n_t$，节点 $i$ 与 $j$ 之间的距离记为 $\ell_{ij}$。对访问顺序 $\sigma=(\sigma_1,\ldots,\sigma_{n_t})$，开放路径长度为

\begin{equation}
L(\sigma)=
\ell_{\mathrm{cur},\sigma_1}
+
\sum_{i=1}^{n_t-1}
\ell_{\sigma_i\sigma_{i+1}},
\label{eq:q4_open_path}
\end{equation}

其中 $\ell_{\mathrm{cur},i}=\lVert P_i-P_{\mathrm{cur}}\rVert_2$ 为当前机器狗位置 $P_{\mathrm{cur}}$ 到节点 $i$ 的欧氏距离，路径不要求返回起点。当前节点集上的调度问题为

\begin{equation}
\min_{\sigma}\quad L(\sigma).
\label{eq:q4_route_problem}
\end{equation}

与问题三相同，$L(\sigma)/5$ 只刻画当前节点集上的移动耗时；执行首节点后，观测结果可能改变非移动动作数和后续节点集，因而算法随证据更新重新规划，并以式~\eqref{eq:q4_total_time} 核算最终虚拟时间。

节点数不超过 10 时，使用 Held--Karp 动态规划求精确开放路径；节点数超过 10 时，先以最小增量插入构造初始解，再使用最近邻起点、2-opt 逆交换和 Or-opt 邻域重定位进行改善。后一种情形只报告高质量可行路径，不宣称全局最优。扫描点集合在规划过程中保持不变，因此重排访问顺序不会改变方向完备发现证书。

\subsubsection{算法流程}

问题四沿用式~\eqref{eq:q4_route_problem}的联合路径求解，但必须改写证据更新与末端清除。下列伪代码中，
$\operatorname{Update}_4$ 对 \texttt{direction} 执行示向锥裁剪、对 \texttt{near} 就地清除，
却不以单次 \texttt{no\_signal} 排除任何位置；除“已发现 16 个异频源”的数量上界外，
$\operatorname{Certify}_4$ 只有在全部方向完备扫描站均无正响应时才能认证该频道不存在。
粗测站包是在首次示向度局部坐标系中对当前外包围的外接矩形作间距不超过 $450\,\mathrm m$ 的有限测站网格，仅用于争取新增正示向度；无信号不收缩定向源外包围，未收敛时仍转入后续带证书的测站包或覆盖清除。

\begin{mdframed}[
  nobreak=true,
  linewidth=0.6pt,
  innerleftmargin=7pt,
  innerrightmargin=7pt,
  innertopmargin=5pt,
  innerbottommargin=2pt,
  skipabove=6pt,
  skipbelow=6pt
]
\textbf{算法 2\quad 问题四方向完备联合调度}
\begin{tabbing}
00\quad\=\quad\quad\=\quad\quad\=\quad\quad\=\kill
1 \> 连续域认证 $P_4'$；若认证失败，则进入 25 点双环工程回退分支并校验上述两项构造不等式；\\
2 \> 调用一次 \texttt{/enter}；初始化频道轨迹 $T_c$、未访问扫描站集 $U$、空清除计数 $b\leftarrow0$；\\
3 \> \textbf{while} 存在未闭合频道 \textbf{do}\\
4 \>\> $\operatorname{Certify}_4(T,U)$；构造扫描节点与服务节点的联合开放路径并取首节点 $a$；\\
5 \>\> \textbf{if} $a$ 为扫描站 $S_i$ \textbf{then}\\
6 \>\>\> 检测必测频道和有严格净收益的顺路频道，逐一执行 $\operatorname{Update}_4$；\\
7 \>\>\> 对仍未知的频道，\texttt{no\_signal} 只更新该站已扫描标记，不收缩 $A_c$；\\
8 \>\> \textbf{else}\quad 令 $c$ 为服务频道，在可见核或保守候选点中选择代价最小测站；\\
9 \>\>\> 最多自适应检测 3 次；遇无信号或收缩不足 $10\%$ 时转入稳健分支；\\
10\>\>\> 若 $\rho_c>300\,\mathrm m$，先执行粗测站包；进入 $300\,\mathrm m$ 内后至多执行 8 轮 12 点包；\\
11\>\>\> 若 $\rho_c\le19.999\,\mathrm m$，则按顶点距离证书清除；\\
12\>\>\> \textbf{else if} $19.5\,\mathrm m$ 覆盖单元所需空清除数不超过 $5-b$ \textbf{then}\\
13\>\>\>\> 按最短顺序覆盖式清除，跳过已排除单元，并以最后单元覆盖证书封口；\\
14\>\>\> \textbf{else} 缩小单元半径直至 12 点包误差界不超过 $19\,\mathrm m$，逐单元测向；\\
15\>\> \textbf{end if}；若覆盖耗尽仍未闭合，则置为 \texttt{geometry\_fault} 并停止；\\
16\> \textbf{end while}；核验 20 个频道证书、$b\le5$ 及清除数范围，仅调用一次 \texttt{/exit}。
\end{tabbing}
\end{mdframed}

\subsection{模型求解与结果分析}

\subsubsection{证书与演练结果}

未提速基础版共完成 6 局演练，平均定位清除时间为 $845.8\mathrm{s}$/源，最好与最坏分别为 $660.7\mathrm{s}$/源和 $1096.9\mathrm{s}$/源，所有实际干扰源均被清除。
完全提速策略定型后共完成 18 局演练，合计清除 $240/240$ 个实际干扰源；逐局平均每源时间的等权均值为 $495.54\mathrm{s}$/源，中位数为 $496.03\mathrm{s}$/源，范围为 $351.95$--$619.72\mathrm{s}$/源，相对基础版平均值降低 $41.41\%$。18 局共发生 12 次空清除，平均每局 $0.67$ 次、单局最多 5 次，均未突破工程预算。
该 18 局跨越算法定型后的两个冻结构建，均采用 22 点连续域认证、覆盖式清除和联合路径调度，故按“完全提速版”统一统计，不再单列中间版本。基础版与完全提速版均保持方向完备与覆盖清除证书；完全提速版 18 局演练及三次正式测试均未触发 25 点回退。

\subsubsection{正式测试结果}

三次正式测试均正常退出，结果如表~\ref{tab:q4_formal}所示。平均定位清除时间和清除比例沿用问题三中 $\bar t_{\mathrm{clr}}$ 与 $\lambda_{\mathrm{clr}}$ 的定义；表中单局平均时间由该局定位清除总时间除以清除个数得到，程序运行时间为模拟器记录的实际运行时间。

\begin{table}[htbp]
\centering
\caption{问题四正式测试结果}
\label{tab:q4_formal}
\small
\setlength{\tabcolsep}{2pt}
\begin{tabular}{@{}clrrr@{}}
\toprule
测试序号 & 测试案例编码 & 清除干扰源个数 & 平均定位清除时间/s & 程序运行时间/ms \\
\midrule
1 & NGBV-4HQZ-7B7S-26HK & 11 & 519.99 & 2426 \\
2 & JGN3-394M-SDZB-Y3N7 & 13 & 493.46 & 2365 \\
3 & 9U2N-ZBZQ-8AB2-VJYS & 15 & 374.53 & 1976 \\
\midrule
合计／总体平均／区间 &  & 39 & 455.20 & 1976--2426 \\
\bottomrule
\end{tabular}
\end{table}



按题面合并口径，每个被清除干扰源的平均定位清除时间为 $17752.83/39=455.20\mathrm{s}$，程序运行时间列给出三次取值区间。三次单局指标的等权算术平均为 $462.66\mathrm{s}$/源，仅作为案例间对照。
三次合计空清除 2 次，均在 5 次预算内；每局 20 个频道最终都已清除或认证为空。
由定理 1，认证为空的频道必无源，因此三次测试均有 $N_{\mathrm{clr}}=N_J$，合并清除比例为 $\lambda_{\mathrm{clr}}=39/39=100\%$，题面要求的两个统计量均已给出。

\subsubsection{时间构成与结果分析}

三次正式测试合计移动距离为 $67.22\mathrm{km}$，时间构成如表~\ref{tab:q4_time}所示。

\begin{table}[htbp]
\centering
\caption{问题四三次正式测试的合计时间构成}
\label{tab:q4_time}
\begin{tabular}{lrr}
\toprule
时间项 & 数值/s & 占比 \\
\midrule
移动时间 & 13443.83 & $75.73\%$ \\
检测时间 & 3470.00 & $19.55\%$ \\
换频时间 & 638.00 & $3.59\%$ \\
成功清除时间 & 195.00 & $1.10\%$ \\
失败清除时间 & 6.00 & $0.03\%$ \\
\midrule
合计 & 17752.83 & $100.00\%$ \\
\bottomrule
\end{tabular}
\end{table}

移动时间占总时间的 $75.73\%$，远高于其他项，说明任务时间主要由路径长度决定；三次移动占比分别为 $73.32\%$、$75.07\%$ 和 $78.92\%$，差异来自干扰源空间分布。
动态开放路径整合扫描站与清除服务点并持续重规划，减少重复往返；22 点扫描保证发现完备性。

三次程序运行时间为 $1.976$--$2.426\mathrm{s}$，远低于 $1200\mathrm{s}$ 上限，不构成约束。
日志复核显示，时间项之和与虚拟时间的最大残差不超过 $2\times10^{-6}\mathrm{s}$；扫描证书、频道状态和动作序列逐项一致，未出现非接受记录或重复站位。

该策略仅使用在线可得的频道号、位置、示向度和清除结果，不使用干扰源、朝向或接收半径真值；发现、定位和盲区清除分别由方向完备证书、外包围和与朝向无关的覆盖保证。
这些结果只说明该策略在题设及问题四的三次正式测试中完成了证书闭合；22 点数和覆盖参数属于工程实现选择，不据此声称对所有随机案例平均最优。
